\documentclass[11pt]{amsart}
\usepackage[margin=1.2in]{geometry}
\usepackage{amssymb,amsmath,amsthm,mathtools}
\usepackage{url}
\usepackage[hidelinks]{hyperref}

% ---------- sync markers (read by scripts/check_sync.py) ----------
% \lean{A.b, A.c}  : Lean declarations realizing this item
% \leanok          : those declarations are fully proved (no sorry, standard axioms)
% \notready        : deliberately not formalized yet
\newcommand{\lean}[1]{\par\smallskip\noindent{\footnotesize\raggedright\textsf{Lean:} \path{#1}\par}}
\newcommand{\leanok}{{\footnotesize\ \textsf{[formalized]}}}
\newcommand{\notready}{\par\smallskip\noindent{\footnotesize\textsf{[not yet formalized]}}}

\newtheorem{theorem}{Theorem}[section]
\newtheorem{proposition}[theorem]{Proposition}
\newtheorem{lemma}[theorem]{Lemma}
\newtheorem{corollary}[theorem]{Corollary}
\theoremstyle{definition}
\newtheorem{definition}[theorem]{Definition}
\theoremstyle{remark}
\newtheorem{remark}[theorem]{Remark}

\DeclareMathOperator{\Aut}{Aut}
\DeclareMathOperator{\ord}{ord}
\newcommand{\Z}{\mathbb{Z}}

\begin{document}
\title{Finite groupoids, correspondences, and divisibility}
\date{}
\begin{abstract}
A short note, kept in sync with the Lean library \texttt{GroupoidCorr}. It formalizes the structural part of the theta-chain account of the generalized Riemann hypothesis. That part has three pieces: integral pushforward along maps of finite groupoids, composition of correspondences, and the fact that divisibility orders add along a chain. Each numbered item names the Lean declarations that realize it. The script \texttt{scripts/check\_sync.py} checks that the two stay in correspondence.
\end{abstract}
\maketitle

Throughout, $R$ is a commutative ring. In the application, $R$ is the ring of integers of a number field, or its localization at a prime.

%======================================================================
\section{Invariants, coinvariants and the norm map}
%======================================================================

Let $G$ be a finite group acting $R$-linearly on an $R$-module $V$ by $\rho$. The \emph{norm} of $v\in V$ is $\sum_{g\in G}\rho(g)v$.

\begin{lemma}\label{lem:norm-map}
\begin{enumerate}
\item If $v$ is invariant, its norm is $|G|\,v$.
\item For every $v$, $|G|\,v=\sum_g\bigl(v-\rho(g)v\bigr)+\sum_g\rho(g)v$.
\item If the norm of $v$ vanishes, then $|G|\,v$ lies in the span of the elements $v-\rho(g)v$.
\end{enumerate}
Consequently the norm map from coinvariants to invariants has kernel and cokernel killed by $|G|$.
\lean{GroupoidCorr.norm_of_invariant, GroupoidCorr.card_smul_eq_norm_add, GroupoidCorr.card_smul_mem_span_of_norm_eq_zero}\leanok
\end{lemma}

\begin{proof}
(1) Every summand equals $v$. (2) Expand $\sum_g(v-\rho(g)v)=|G|\,v-\sum_g\rho(g)v$. (3) Substitute the vanishing norm into (2).

For the consequence: by (1), the image of the norm map contains $|G|$ times every invariant, so $|G|$ kills the cokernel. By (3), $|G|$ kills the kernel, since the span of the $v-\rho(g)v$ is the kernel of the projection to coinvariants.
\end{proof}

%======================================================================
\section{Local systems on groupoids and pushforward}
%======================================================================

\begin{definition}\label{def:local-system}
A \emph{local system} of $R$-modules on a groupoid $X$ is a functor $F\colon X\to R\text{-Mod}$.
\lean{GroupoidCorr.LocalSystem}\leanok
\end{definition}

\begin{definition}\label{def:section}
A \emph{section} of $F$ is a family $(s_x)_{x\in X}$ with $s_x\in F(x)$ such that $F(g)(s_x)=s_y$ for every morphism $g\colon x\to y$.
\lean{GroupoidCorr.Sec}\leanok
\end{definition}

\begin{definition}\label{def:fibre}
Let $f\colon X\to Y$ be a functor of groupoids and $y\in Y$. The \emph{fibre} of $f$ over $y$ is the set of equivalence classes of pairs $(x,\alpha)$ with $\alpha\colon f(x)\to y$. Two pairs are equivalent, $(x,\alpha)\sim(x',\alpha')$, if some $\beta\colon x\to x'$ satisfies $f(\beta);\alpha'=\alpha$. Here $;$ denotes composition in diagrammatic order (first $f(\beta)$, then $\alpha'$), matching the composition operator of Lean's category theory library.
\lean{GroupoidCorr.fibreSetoid, GroupoidCorr.Fibre}\leanok
\end{definition}

\begin{proof}[The relation is an equivalence relation]
Reflexivity: take $\beta=\mathrm{id}$. Symmetry: if $f(\beta);\alpha'=\alpha$, then $f(\beta^{-1});\alpha=f(\beta^{-1});f(\beta);\alpha'=\alpha'$. Transitivity: if $f(\beta);\alpha'=\alpha$ and $f(\gamma);\alpha''=\alpha'$, then $f(\beta;\gamma);\alpha''=f(\beta);f(\gamma);\alpha''=\alpha$.
\end{proof}

\begin{lemma}\label{lem:summand-wd}
Let $F$ be a local system on $Y$ and $s$ a section of $f^*F=F\circ f$. The element $F(\alpha)(s_x)\in F(y)$ depends only on the class of $(x,\alpha)$ in the fibre.
\lean{GroupoidCorr.summand_wd}\leanok
\end{lemma}

\begin{proof}
If $f(\beta);\alpha'=\alpha$, then
\[
F(\alpha')(s_{x'})=F(\alpha')\bigl(F(f\beta)(s_x)\bigr)=F\bigl(f(\beta);\alpha'\bigr)(s_x)=F(\alpha)(s_x),
\]
using the compatibility of $s$ along $\beta$ and the functoriality of $F$.
\end{proof}

\begin{definition}\label{def:pushforward}
Assume all fibres of $f$ are finite. The \emph{pushforward} of a section $s$ of $f^*F$ is
\[
(f_!s)_y=\sum_{[(x,\alpha)]\in\text{fibre over }y}F(\alpha)(s_x)\ \in F(y).
\]
It is a finite sum with integer coefficients: nothing is divided.
\lean{GroupoidCorr.pushVal}\leanok
\end{definition}

\begin{proposition}\label{prop:push-section}
$f_!s$ is a section of $F$.
\lean{GroupoidCorr.pushVal_compat, GroupoidCorr.push}\leanok
\end{proposition}

\begin{proof}
Let $g\colon y\to y'$. Post-composition $(x,\alpha)\mapsto(x,\alpha;g)$ respects the equivalence relation (if $f(\beta);\alpha'=\alpha$ then $f(\beta);\alpha';g=\alpha;g$), and it is a bijection from the fibre over $y$ to the fibre over $y'$, with inverse given by post-composition with $g^{-1}$. It carries the summand $F(\alpha)(s_x)$ to $F(\alpha;g)(s_x)=F(g)\bigl(F(\alpha)(s_x)\bigr)$. Since $F(g)$ is additive, summing gives $F(g)\bigl((f_!s)_y\bigr)=(f_!s)_{y'}$.
\end{proof}

\begin{definition}\label{def:pull-map}
Let $f\colon X\to Y$ be a functor and $s$ a section of a local system $F$ on $Y$. The \emph{pullback} $f^*s$ is the section $(s_{f(x)})_x$ of $f^*F$. If $K\colon F\to G$ is a morphism of local systems on $X$, that is, a natural transformation, then $K(s)=(K_x(s_x))_x$ is a section of $G$.
\lean{GroupoidCorr.Sec.pull, GroupoidCorr.Sec.map}\leanok
\end{definition}

\begin{lemma}\label{lem:push-natural}
Let $k\colon A\to B$ have finite fibres, $K\colon G\to G'$ a morphism of local systems on $B$, and $s$ a section of $k^*G$. Then $K_b\bigl((k_!s)_b\bigr)=\bigl(k_!(k^*K)(s)\bigr)_b$ for every $b$.
\lean{GroupoidCorr.pushVal_map}\leanok
\end{lemma}

\begin{proof}
$K_b$ is additive, and $K_b\bigl(G(\beta)(s_a)\bigr)=G'(\beta)\bigl(K_{k(a)}(s_a)\bigr)$ by naturality.
\end{proof}

\begin{lemma}\label{lem:push-comp}
Let $k\colon A\to B$ and $h\colon B\to C$ have finite fibres, with $h$ \emph{faithful}. Then the fibres of $h\circ k$ are finite, and $(h\circ k)_!s=h_!(k_!s)$ for every section $s$ of $(h\circ k)^*F$.
\lean{GroupoidCorr.pushVal_comp}\leanok
\end{lemma}

\begin{proof}
Fix $c\in C$ and choose a representative $(b_q,\alpha_q)$ of each class $q$ in the fibre of $h$ over $c$. The map
\[
\textstyle\coprod_q\,(\text{fibre of }k\text{ over }b_q)\longrightarrow\text{fibre of }h\circ k\text{ over }c,\qquad (q,[(a,\beta)])\mapsto[(a,h(\beta);\alpha_q)]
\]
is well defined and surjective: given $(a,\varepsilon)$, let $q=[(k(a),\varepsilon)]$ and pick $\beta\colon k(a)\to b_q$ with $h(\beta);\alpha_q=\varepsilon$. This proves finiteness. It is injective when $h$ is faithful. If $(a,h\beta;\alpha_q)\sim(a',h\beta';\alpha_{q'})$ via $\gamma$, then $\beta^{-1};k(\gamma);\beta'$ shows $q=q'$. Cancelling $\alpha_q$ and using faithfulness gives $k(\gamma);\beta'=\beta$. The summand at $(q,[(a,\beta)])$ is $F(h\beta;\alpha_q)(s_a)=F(\alpha_q)\bigl(F(h\beta)(s_a)\bigr)$. Summing over the inner fibre gives the summand $F(\alpha_q)\bigl((k_!s)_{b_q}\bigr)$ of $h_!(k_!s)$.
\end{proof}

\begin{remark}
Faithfulness of $h$ cannot be dropped. For a finite group $G$, take $k\colon\mathrm{pt}\to BG$, $h\colon BG\to\mathrm{pt}$ and a constant local system. Then $h_!k_!$ is multiplication by $|G|$, while $(h\circ k)_!$ is the identity. This is why the main account requires the target legs of correspondences to be faithful. Pushforward itself, and base change, need no faithfulness.
\end{remark}

\begin{definition}\label{def:fibre-product}
Let $f\colon X\to Y$ and $g\colon Z\to Y$ be functors of groupoids. The \emph{$2$-fibre product} $W=X\times_YZ$ has as objects the triples $(x,z,\gamma)$ with $\gamma\colon f(x)\to g(z)$. A morphism $(x,z,\gamma)\to(x',z',\gamma')$ is a pair $(\beta\colon x\to x',\ \delta\colon z\to z')$ with $f(\beta);\gamma'=\gamma;g(\delta)$. It is a groupoid, with $(\beta,\delta)^{-1}=(\beta^{-1},\delta^{-1})$. Write $g'\colon W\to X$ and $f'\colon W\to Z$ for the projections. If $f$ is faithful, so is $f'$: if $(\beta,\delta)$ and $(\beta_1,\delta)$ are morphisms $(x,z,\gamma)\to(x',z',\gamma')$, then $f(\beta);\gamma'=\gamma;g(\delta)=f(\beta_1);\gamma'$, so $f(\beta)=f(\beta_1)$ and $\beta=\beta_1$. In Lean, $W$ is Mathlib's comma category $\mathrm{Comma}(f,g)$. Mathlib does not provide its groupoid structure, so we build it by hand.
\lean{GroupoidCorr.commaGroupoid, GroupoidCorr.Corr.snd_faithful}\leanok
\end{definition}

\begin{proposition}[base change]\label{prop:base-change}
In the setting of Definition~\ref{def:fibre-product}, assume the fibres of $f$ are finite. Let $F$ be a local system on $Y$ and $s$ a section of $f^*F$. Then:
\begin{enumerate}
\item The family $\tilde s_{(x,z,\gamma)}=F(\gamma)(s_x)$ is a section of $f'^*g^*F$. This is $g'^*s$, transported along the natural isomorphism $\gamma\colon g'^*f^*F\cong f'^*g^*F$.
\item For each $z\in Z$, post-composition with $g$ gives a bijection from the fibre of $f'$ over $z$ to the fibre of $f$ over $g(z)$. In particular, the fibres of $f'$ are finite.
\item $(f'_!\tilde s)_z=(f_!s)_{g(z)}$ for every $z\in Z$. That is, $g^*f_!=f'_!\,g'^*$.
\end{enumerate}
\lean{GroupoidCorr.baseChangeSec, GroupoidCorr.bcFibreEquiv, GroupoidCorr.pushVal_baseChange}\leanok
\end{proposition}

\begin{proof}
(1) For a morphism $(\beta,\delta)$, functoriality and the compatibility of $s$ along $\beta$ give
\[
F(g\delta)\bigl(F(\gamma)(s_x)\bigr)=F(\gamma;g\delta)(s_x)=F\bigl(f\beta;\gamma'\bigr)(s_x)=F(\gamma')(s_{x'}).
\]
(2) The map sends the class of $((x,z_0,\gamma),\delta\colon z_0\to z)$ to the class of $(x,\gamma;g(\delta))$. A morphism $(\beta,\delta_1)$ between pairs maps to $\beta$. The inverse sends $(x,\alpha)$ to $((x,z,\alpha),\mathrm{id}_z)$. The composite on the fibre of $f'$ is related to the identity by the morphism $(\mathrm{id}_x,\delta^{-1})$. The other composite is the identity.
(3) The bijection of (2) carries the summand $F(g\delta)\bigl(F(\gamma)(s_x)\bigr)$ to $F(\gamma;g\delta)(s_x)$. These are equal by functoriality. Sum over the fibre.
\end{proof}

\begin{definition}\label{def:finite-groupoid}
A groupoid is \emph{finite} if it has finitely many isomorphism classes of objects and every automorphism group is finite. Then every hom-set is finite, and the fibres of any functor between finite groupoids are finite. A $\mathbb C$-valued function on isomorphism classes is the same as a section of the constant local system $\mathbb C$. For two such functions put
\[
\langle a,b\rangle_X=\sum_{[x]}\frac{a(x)\,\overline{b(x)}}{|\Aut x|},
\]
computed on chosen representatives of the classes. Because $a$, $b$ and $|\Aut x|$ are constant on isomorphism classes, the result does not depend on the choice.
\lean{GroupoidCorr.FinGroupoid, GroupoidCorr.ginner}\leanok
\end{definition}

\begin{proposition}[groupoid cardinality and adjointness]\label{prop:adjoint}
Let $f\colon X\to Y$ be a faithful functor of finite groupoids, $a$ a section of $f^*\mathbb C$ and $b$ a section of $\mathbb C$ on $Y$. Then $\langle f_!a,b\rangle_Y=\langle a,f^*b\rangle_X$.
\lean{GroupoidCorr.ginner_push}\leanok
\end{proposition}

\begin{proof}
Fix representatives $x$ of the classes of $X$ and $y$ of those of $Y$. Count the triples $([x],[y],\alpha\colon f(x)\to y)$ in two ways.

First fix $y$. The map $([x],\alpha)\mapsto[(x,\alpha)]$ onto the fibre over $y$ is surjective. The triples over the class $[(x',\alpha')]$ are the $([x'],f(\beta);\alpha')$ with $\beta\colon x\to x'$, where $x$ is the representative of the class of $x'$. These are pairwise distinct by faithfulness, so there are $|\Aut x'|$ of them. Hence
\[
(f_!a)(y)=\sum_{[x]}\sum_{\alpha\colon f(x)\to y}\frac{a(x)}{|\Aut x|}.
\]
Now fix $x$. The set $\mathrm{Hom}(f(x),y)$ is empty unless $[y]=[f(x)]$, and then it has $|\Aut y|$ elements. Therefore
\[
\langle f_!a,b\rangle_Y=\sum_{[x]}\sum_{[y]}\bigl|\mathrm{Hom}(f(x),y)\bigr|\,\frac{a(x)}{|\Aut x|}\,\frac{\overline{b(y)}}{|\Aut y|}=\sum_{[x]}\frac{a(x)\,\overline{b(f(x))}}{|\Aut x|}=\langle a,f^*b\rangle_X.
\]
This is the orbit--stabilizer formula for the action of $\Aut(y)$ on the fibre over $y$, written as a count.
\end{proof}

%======================================================================
\section{Correspondences, composition, and divisibility}
%======================================================================

\begin{definition}\label{def:correspondence}
A \emph{correspondence} from $(X,F)$ to $(X',F')$ consists of:
\begin{itemize}
\item a groupoid $C$ (in Lean, a parameter of the structure);
\item functors $\mathrm{src}\colon C\to X$ and $\mathrm{tgt}\colon C\to X'$, with $\mathrm{tgt}$ faithful and with finite fibres;
\item a \emph{kernel} $K\colon\mathrm{src}^*F\to\mathrm{tgt}^*F'$, a morphism of local systems on $C$.
\end{itemize}
Its operator, on sections, is $T_C=\mathrm{tgt}_!\circ K\circ\mathrm{src}^*$.
\lean{GroupoidCorr.Corr, GroupoidCorr.Corr.op}\leanok
\end{definition}

\begin{proposition}[composition]\label{prop:composition}
Let $C_1$ be a correspondence from $(X,F)$ to $(X',F')$ and $C_2$ one from $(X',F')$ to $(X'',F'')$. Let $W=C_1\times_{X'}C_2$ be the $2$-fibre product of $\mathrm{tgt}_1$ and $\mathrm{src}_2$, with projections $g'\colon W\to C_1$ and $f'\colon W\to C_2$. Give $W$ the legs $\mathrm{src}_1\circ g'$ and $\mathrm{tgt}_2\circ f'$, and the kernel
\[
K_{(c_1,c_2,\gamma)}=K_{2,c_2}\circ F'(\gamma)\circ K_{1,c_1}.
\]
This is a correspondence $C_2\circ C_1$ from $(X,F)$ to $(X'',F'')$, and $T_{C_2}\circ T_{C_1}=T_{C_2\circ C_1}$.
\lean{GroupoidCorr.Corr.comp, GroupoidCorr.Corr.op_comp}\leanok
\end{proposition}

\begin{proof}
The kernel is natural: for a morphism $(\beta,\delta)$ of $W$, use the naturality of $K_1$ along $\beta$, then $\mathrm{tgt}_1(\beta);\gamma'=\gamma;\mathrm{src}_2(\delta)$, then the naturality of $K_2$ along $\delta$. The leg $\mathrm{tgt}_2\circ f'$ is faithful by Definition~\ref{def:fibre-product}. Its fibres are finite by Proposition~\ref{prop:base-change}(2) and Lemma~\ref{lem:push-comp}. For the identity, write $u=K_1(\mathrm{src}_1^*s)$. Then
\[
K_2\,\mathrm{src}_2^*\,\mathrm{tgt}_{1!}u=K_2\,f'_!\tilde u=f'_!\bigl((f'^*K_2)(\tilde u)\bigr)
\]
by base change (Proposition~\ref{prop:base-change}) and Lemma~\ref{lem:push-natural}. Apply $\mathrm{tgt}_{2!}$ and use Lemma~\ref{lem:push-comp} for the faithful $\mathrm{tgt}_2$. This gives $(\mathrm{tgt}_2\circ f')_!$ of the section $K_{2,c_2}\bigl(F'(\gamma)(K_{1,c_1}(s_{\mathrm{src}_1c_1}))\bigr)$, which is $T_{C_2\circ C_1}s$.
\end{proof}

Operators of correspondences are $R$-linear, since pullback, kernels and pushforward are. So the divisibility statement reduces to linear algebra.

\begin{definition}\label{def:chain}
For $R$-linear maps $K_i\colon M_i\to M_{i+1}$, the composite of the first $n$ steps is
\[
K_{(n)}=K_{n-1}\circ\cdots\circ K_0\colon M_0\to M_n,\qquad K_{(0)}=\mathrm{id}.
\]
\lean{GroupoidCorr.chainComp}\leanok
\end{definition}

\begin{theorem}[orders add along a chain]\label{thm:orders-add}
Let $\varpi\in R$, and suppose $K_i=\varpi^{r_i}K_i'$ for $R$-linear maps $K_i'$. Then
\[
K_{(n)}=\varpi^{\,r_0+\cdots+r_{n-1}}\,K'_{(n)} .
\]
\lean{GroupoidCorr.chainComp_eq_pow_smul}\leanok
\end{theorem}

\begin{proof}
Induction on $n$. For $n=0$ both sides are the identity. For the step,
\[
K_{(n+1)}=K_n\circ K_{(n)}=\bigl(\varpi^{r_n}K_n'\bigr)\circ\bigl(\varpi^{S}K'_{(n)}\bigr)=\varpi^{S+r_n}\,K'_{(n+1)},
\]
where $S=r_0+\cdots+r_{n-1}$, because composition is bilinear.
\end{proof}

\begin{corollary}\label{cor:divisible-values}
If a bilinear period pairing satisfies $\mathit{per}=\varpi^{S}\mathit{per}'$ with $\mathit{per}'$ integral, then every value of $\mathit{per}$ lies in the ideal $(\varpi^{S})$.
\lean{GroupoidCorr.values_mem_span_pow}\leanok
\end{corollary}

\begin{proof}
$\mathit{per}(x)(y)=\varpi^{S}\cdot\mathit{per}'(x)(y)$.
\end{proof}

The same statement holds one level up, for the kernels themselves.

\begin{lemma}\label{lem:corr-smul}
For $c\in R$ and a correspondence $C$, write $cC$ for the same groupoid and legs with kernel $cK$. Then $T_{cC}=c\,T_C$, and $(dC_2)\circ(cC_1)=cd\,(C_2\circ C_1)$.
\lean{GroupoidCorr.Corr.smul, GroupoidCorr.Corr.op_smul, GroupoidCorr.Corr.comp_smul}\leanok
\end{lemma}

\begin{proof}
Pullback and pushforward are $R$-linear, so the first identity holds summand by summand. For the second, the composite kernel $K_{2,c_2}\circ F'(\gamma)\circ K_{1,c_1}$ is bilinear in $(K_1,K_2)$.
\end{proof}

\begin{definition}\label{def:corr-id}
The \emph{identity correspondence} on $(X,F)$ has groupoid $X$, both legs the identity and kernel the identity. Its fibres are single points, and its operator is the identity.
\lean{GroupoidCorr.Corr.id, GroupoidCorr.Corr.op_id}\leanok
\end{definition}

\begin{definition}\label{def:corr-chain}
Given correspondences $C_i$ from $(X_i,F_i)$ to $(X_{i+1},F_{i+1})$, put $C_{(0)}$ equal to the identity correspondence on $(X_0,F_0)$, and $C_{(n+1)}=C_n\circ C_{(n)}$ (Proposition~\ref{prop:composition}). Then $T_{C_{(0)}}=\mathrm{id}$ and $T_{C_{(n+1)}}=T_{C_n}\circ T_{C_{(n)}}$. In Lean the groupoid of a correspondence is a field (\texttt{BCorr}), because it changes along the chain. All groupoids are taken in a single universe, which the $2$-fibre product preserves.
\lean{GroupoidCorr.chain, GroupoidCorr.chain_op_zero, GroupoidCorr.chain_op_succ}\leanok
\end{definition}

\begin{theorem}[orders add along a chain of correspondences]\label{thm:orders-add-corr}
Let $\varpi\in R$, and suppose $C_i=\varpi^{r_i}C_i'$ for correspondences $C_i'$, in the sense of Lemma~\ref{lem:corr-smul}. Then
\[
C_{(n)}=\varpi^{\,r_0+\cdots+r_{n-1}}\,C'_{(n)},\qquad\text{hence}\qquad T_{C_{(n)}}=\varpi^{\,r_0+\cdots+r_{n-1}}\,T_{C'_{(n)}}.
\]
\lean{GroupoidCorr.chain_eq_pow_smul, GroupoidCorr.chain_op_eq_pow_smul}\leanok
\end{theorem}

\begin{proof}
Induction on $n$, as in Theorem~\ref{thm:orders-add}, with Lemma~\ref{lem:corr-smul} in place of the bilinearity of composition. For $n=0$, $1\cdot C=C$. The operator statement follows from $T_{cC}=c\,T_C$.
\end{proof}

\begin{remark}
Theorem~\ref{thm:orders-add-corr}, with Corollary~\ref{cor:divisible-values} applied to the period pairing, gives the divisibility of the periods of a theta chain once each step's kernel is divisible by $\varpi^{r_i}$. This is Theorem~8.11 and Corollary~8.12 of the main account. Theorem~\ref{thm:orders-add} is the same statement for the operators alone, which are $R$-linear maps.
\end{remark}

%======================================================================
\section{Denominators of the Lubin--Tate exponential}
%======================================================================

This section is Section~8.2 of the main account. There $\mathcal O$ is a complete discrete valuation ring, unramified over $\Z_p$, with residue field $\mathbb F_q$. $F$ is a formal group over $\mathcal O$ with $[p]_F(T)=pT+\sum_{k\ge2}e_kT^k\equiv T^q\pmod p$, and $\exp_F(Z)=\sum_{n\ge1}c_nZ^n$ is the inverse of its normalized logarithm. The formal group enters only through the identity $[p]_F(\exp_F Z)=\exp_F(pZ)$, so we take that identity as a hypothesis.

To avoid fractional valuations, we adjoin a root $\varpi$ of $\varpi^{q-1}=p$. Then $\ord_p(x)\ge-m/(q-1)$ means exactly that $\varpi^{m}x$ is integral. The only consequence of ``unramified'' that is used is that $p^m-1$ is a unit for $m\ge1$.

\begin{theorem}[denominators of the Lubin--Tate exponential]\label{thm:LT-exp}
Let $O$ be a subring of a field $K$, and let $q\ge2$. Let $\varpi\in O$ with $\varpi^{q-1}=p\ne0$, and suppose $p^m-1$ is a unit of $O$ for every $m\ge1$. Let $e_k\in O$ for $k\ge2$, with $e_k\in pO$ for $k\ne q$. Let $E=\sum_nc_nZ^n\in K[[Z]]$ with $c_0=0$ and $c_1=1$, such that for every $n$
\[
p^nc_n=p\,c_n+\sum_{k=2}^{n}e_k\,[Z^n]E^k .
\]
This is the coefficient of $Z^n$ in $[p](E(Z))=E(pZ)$. Then $\varpi^{n-1}c_n\in O$ for all $n\ge1$, that is, $\ord_p(c_n)\ge-(n-1)/(q-1)$.
\lean{GroupoidCorr.LTData, GroupoidCorr.LTData.exp_coeff_mem}\leanok
\end{theorem}

\begin{proof}
Put $G(Z)=\varpi^{-1}E(\varpi Z)=\sum_n g_nZ^n$, so $g_n=\varpi^{n-1}c_n$ and $[Z^n]E^k=\varpi^{k-n}[Z^n]G^k$. We show $g_n\in O$ by strong induction on $n$. For $n\le1$, $g_0=0$ and $g_1=1$. For $n\ge2$, multiply the identity by $\varpi^{n-1}$:
\[
p\,(p^{n-1}-1)\,g_n=\sum_{k=2}^{n}e_k\varpi^{k-1}\,[Z^n]G^k .
\]
For $k\ge2$, $[Z^n]G^k$ involves only $g_m$ with $m<n$, because $g_0=0$. So it lies in $O$ by induction. Each $e_k\varpi^{k-1}$ is divisible by $p$: for $k=q$ because $\varpi^{q-1}=p$, and for $k\ne q$ because $p\mid e_k$. Divide by $p$ and by the unit $p^{n-1}-1$.
\end{proof}

\begin{theorem}[divisibility of jets]\label{thm:jet-gain}
In the setting of Theorem~\ref{thm:LT-exp}, let $f=\sum_mf_mT^m$ with $f_m\in O$, and let $b_a=[Z^a]f(E(Z))=\sum_{m=0}^{a}f_m[Z^a]E^m$. Then $\varpi^{a-1}b_a\in O$ for $a\ge1$. Hence $\varpi^{a-1}\,a!\,b_a\in a!\,O$.

When $\partial$ is the invariant derivation with $\partial\lambda_F=1$, we have $(\partial^af)(0)=a!\,b_a$, since $\partial=d/dZ$ in the coordinate $Z=\lambda_F(T)$. So this says
\[
\ord_p\bigl((\partial^af)(0)\bigr)\ \ge\ \ord_p(a!)-\frac{a-1}{q-1}.
\]
\lean{GroupoidCorr.LTData.jet_coeff_mem, GroupoidCorr.LTData.jet_factorial_mem}\leanok
\end{theorem}

\begin{proof}
The term $m=0$ vanishes because $a\ge1$. For $m\ge1$, $\varpi^{a-1}f_m[Z^a]E^m=f_m\varpi^{m-1}[Z^a]G^m$, and $G$ has integral coefficients by the proof of Theorem~\ref{thm:LT-exp}.
\end{proof}

\begin{corollary}[height decides the gain]\label{cor:height}
Let $p$ be prime, $a\ge1$, and $s=s_p(a)$ the sum of the base-$p$ digits of $a$.
\begin{enumerate}
\item Split primes ($q=p$): $(p-1)\bigl(\ord_p(a!)-\tfrac{a-1}{p-1}\bigr)=1-s\le0$. There is no gain.
\item Inert primes ($q=p^2$): $(p^2-1)\bigl(\ord_p(a!)-\tfrac{a-1}{p^2-1}\bigr)=ap-(p+1)s+1$.
\end{enumerate}
\lean{GroupoidCorr.height_split, GroupoidCorr.height_inert}\leanok
\end{corollary}

\begin{proof}
Legendre's formula gives $(p-1)\ord_p(a!)=a-s$. For (1), subtract $a-1$; also $s\ge1$, since the leading digit is nonzero. For (2), multiply Legendre's formula by $p+1$ and subtract $a-1$.
\end{proof}

\begin{remark}
In case (2) the gain is $\frac{ap}{p^2-1}-\frac{(p+1)s-1}{p^2-1}=\sum_{k\text{ odd}}a/p^k-O(\log_pa)$, since $s\le(p-1)(\log_pa+1)$. This is the pattern $k/p+k/p^3+\cdots-O(\log k)$ for jets of CM theta functions at inert primes (main account, Corollary~8.5).
\end{remark}

\end{document}
